\begin{trapbox}
    \par\noindent\begin{minipage}{\linewidth}
    \textbf{\textcolor{CTrap}{易错一：把竖直中线 $PM$ 当成底边 $AB$ 的高}}

    \begin{center}
\begin{tikzpicture}[x=1.05cm,y=1.05cm,font=\small,>=Stealth]
\path (-2,-2.1) rectangle (9,3.6);
\coordinate (A) at (-1,.25); \coordinate (B) at (3,2.25);
\coordinate (P) at (1,-.75); \coordinate (M) at (1,1.25);
\coordinate (N) at (1,.25); \coordinate (H) at (.2,.85);
\fill[blue!12] (P)--(A)--(M)--cycle;
\fill[orange!20] (P)--(M)--(B)--cycle;
\draw[->,gray] (-1.8,0)--(3.8,0) node[right] {$x$};
\draw[->,gray] (0,-1.35)--(0,3.2) node[above] {$y$};
\draw[blue!70!black,thick,domain=-1.6:3.45,samples=100] plot (\x,{\x*\x/4});
\draw[gray!80,domain=-1.7:1.5] plot (\x,{-.5*\x-.25});
\draw[gray!80,domain=.7:3.4] plot (\x,{1.5*\x-2.25});
\draw[very thick] (A)--(B);
\draw[thick] (A)--(P)--(B);
\draw[orange!80!black,very thick] (P)--(M);
\foreach \pt/\pos in {B/above left,M/above left}
  \fill (\pt) circle (1.6pt) node[\pos] {$\pt$};
\fill (A) circle (1.6pt); \node[anchor=east] at (-1.3,-.35) {$A$};
\fill (P) circle (1.6pt); \node[anchor=west] at (1.4,-.65) {$P$};

\draw[teal!80!black,very thick] (P)--(H);
\fill[teal!80!black] (H) circle (1.6pt) node[above] {$H$};
\draw[teal!80!black] ($(H)+(.161,.0805)$)--($(H)+(.2415,-.0805)$)--($(H)+(.0805,-.161)$);

\node[anchor=north west,align=left,text width=4.7cm] at (4.1,3.1) {\textbf{以 $AB$ 为底，须配 $PH$}\\[8pt]
$PH\perp AB,\quad d=PH$\\[8pt]
$S=\dfrac12 |AB|\,d$\\[10pt]
橙色 $PM$ 经过中点，\\绿色 $PH$ 垂直于底边。\\[8pt]
本图中 $H\ne M$，\\不能以 $PM$ 代替 $d$。};
\end{tikzpicture}
\end{center}

    \end{minipage}\par

    \textbf{错法：} $PM$ 竖直，于是直接写
    $S=\frac12|AB|\cdot PM$。
    \textbf{错误原因：}中线只要求经过边的中点，\textbf{不要求垂直于边}。
    以 $A(2,1),B(4,4),P(3,2)$ 为反例：
    $AB$ 斜率为 $3/2$，故竖直的 $PM$ 不与 $AB$ 垂直。
    \textbf{正确做法：}以\textbf{竖直中线 $PM$} 为底，
    点 $A,B$ 到 $PM$ 的水平距离相加等于 $|x_1-x_2|$，所以
    \[
        \boxed{S=\frac12 PM\cdot|x_1-x_2|}.
    \]
    只有 $x_1+x_2=0$ 时，$AB$ 水平，此时 $PM$ 才同时是 $AB$ 上的高。

    \medskip
    \textbf{\textcolor{CTrap}{易错二：误记系数和绝对值}}

    \textbf{错法：}写成 $S=|x_1-x_2|^3/(8p^2)$ 或
    $S=(x_1-x_2)^3/(8p)$。
    \textbf{正确做法：}
    \[
        \boxed{S=\frac{|x_1-x_2|^3}{8p}},\qquad p>0.
    \]
    例如 $x_1=0,x_2=2,p=1$，面积应为 $8/8=1$；
    不加绝对值并取相反次序 $x_1=0,x_2=2$，直接算
    $(0-2)^3/8=-1$，出现“不可能的负面积”。
    \textbf{原因：}面积为非负量，且横坐标差的次序可互换。

    \medskip
    \par\noindent\begin{minipage}{\linewidth}
    \textbf{\textcolor{CTrap}{易错三：把顶点纵坐标写成横坐标的平方}}

    \begin{center}
\begin{tikzpicture}[x=1.05cm,y=1.05cm,font=\small,>=Stealth]
\path (-2,-2.1) rectangle (9,3.6);
\coordinate (A) at (-1,.25); \coordinate (B) at (3,2.25);
\coordinate (P) at (1,-.75); \coordinate (M) at (1,1.25);
\coordinate (N) at (1,.25); \coordinate (H) at (.2,.85);
\fill[blue!12] (P)--(A)--(M)--cycle;
\fill[orange!20] (P)--(M)--(B)--cycle;
\draw[->,gray] (-1.8,0)--(3.8,0) node[right] {$x$};
\draw[->,gray] (0,-1.35)--(0,3.2) node[above] {$y$};
\draw[blue!70!black,thick,domain=-1.6:3.45,samples=100] plot (\x,{\x*\x/4});
\draw[gray!80,domain=-1.7:1.5] plot (\x,{-.5*\x-.25});
\draw[gray!80,domain=.7:3.4] plot (\x,{1.5*\x-2.25});
\draw[very thick] (A)--(B);
\draw[thick] (A)--(P)--(B);
\draw[orange!80!black,very thick] (P)--(M);
\foreach \pt/\pos in {B/above left,M/above left}
  \fill (\pt) circle (1.6pt) node[\pos] {$\pt$};
\fill (A) circle (1.6pt); \node[anchor=east] at (-1.3,-.35) {$A$};
\fill (P) circle (1.6pt); \node[anchor=west] at (1.4,-.65) {$P$};

\fill[blue!70!black] (N) circle (1.7pt) node[right] {$N$};
\draw[thick] (.88,-.25)--(1.12,-.25) (.88,.75)--(1.12,.75);

\node[anchor=north west,align=left,text width=4.7cm] at (4.1,3.1) {\textbf{平分的是线段 $PM$}\\[8pt]
$x_N=x_P=x_M$\\[8pt]
$y_N=\dfrac{(x_1+x_2)^2}{8p}$\\[10pt]
$PN=NM=\dfrac{PM}{2}$\\[8pt]
$P$ 是切线交点；\\$N$ 才是同横坐标的曲线点。\\[8pt]
两段曲线弧不据此等长。};
\end{tikzpicture}
\end{center}

    \end{minipage}\par

    \textbf{错法：}误以为 $P$ 在抛物线上，写
    $y_P=x_P^2/(2p)$。
    \textbf{正确做法：}
    \[
        x_P=\frac{x_1+x_2}{2},\qquad
        y_P=\frac{x_1x_2}{2p}.
    \]
    当 $x_1\ne x_2$ 时，
    \[
        \frac{x_P^2}{2p}-y_P
        =\frac{(x_1-x_2)^2}{8p}>0,
    \]
    即 $P$ 严格位于抛物线的下方，\textbf{并不在抛物线上}。
    这也解释了为什么不能直接把 $P$ 当成第三个曲线上点。

    \medskip
    \textbf{\textcolor{CTrap}{易错四：两切点重合仍套用三角形}}

    \textbf{错法：}没有确认 $x_1\ne x_2$ 就在“相减后的方程”中
    除以 $x_1-x_2$。
    \textbf{正确做法：}题设必须是两个不同切点。
    如果 $x_1=x_2$，两条所谓切线其实是同一条切线，
    既不能通过两条不同直线确定唯一交点，三角形也不成立。
    已知 $P(x_0,y_0)$ 反求切点时，应先检查
    \[
        x_0^2-2py_0>0.
    \]
    等于零时只有一个切点；小于零时不存在实切点。

    \medskip
    \par\noindent\begin{minipage}{\linewidth}
    \textbf{\textcolor{CTrap}{易错五：忽视开口方向}}

    \begin{center}
\begin{tikzpicture}[x=.9cm,y=.9cm,font=\small,>=Stealth]
\foreach \shift/\swap/\title in {0/0/{$x^2=4y$：竖直中线},7/1/{$y^2=4x$：水平中线}} {
\begin{scope}[xshift=\shift cm]
\node[anchor=south] at (2,5) {\title};
\draw[->,gray] (-.4,0)--(4.9,0) node[right] {$x$};
\draw[->,gray] (0,-.4)--(0,4.7) node[above] {$y$};
\ifnum\swap=0
\coordinate (A) at (2,1); \coordinate (B) at (4,4);
\coordinate (P) at (3,2); \coordinate (M) at (3,2.5);
\draw[blue!70!black,thick,domain=0:4.2,samples=70] plot (\x,{\x*\x/4});
\node at (3.5,1.8) {$P$}; \node at (2.5,2.7) {$M$};
\else
\coordinate (A) at (1,2); \coordinate (B) at (4,4);
\coordinate (P) at (2,3); \coordinate (M) at (2.5,3);
\draw[blue!70!black,thick,domain=0:4.2,samples=70] plot ({\x*\x/4},\x);
\node at (1.6,3.4) {$P$}; \node at (2.8,2.7) {$M$};
\fi
\fill[blue!12] (P)--(A)--(B)--cycle;
\draw[thick] (P)--(A)--(B)--cycle;
\draw[orange!80!black,very thick] (P)--(M);
\foreach \pt/\pos in {A/below right,B/above left}
  \fill (\pt) circle (1.5pt) node[\pos] {$\pt$};
\fill (P) circle (1.5pt); \fill (M) circle (1.5pt);
\end{scope}}
\end{tikzpicture}
\end{center}

    \end{minipage}\par

    \textbf{错法：}把 $x^2=2py$ 的公式原封不动用于 $y^2=2px$。
    \textbf{正确做法：}坐标交换后应写为
    \[
        P\left(\frac{y_1y_2}{2p},\frac{y_1+y_2}{2}\right),\qquad
        S=\frac{|y_1-y_2|^3}{8p}.
    \]
    这时是\textbf{水平}中线，而不是竖直中线。

    \medskip
    \textbf{\textcolor{CTrap}{易错六：把“中线被平分”理解成“弧长被平分”}}

    令竖直线 $x=x_P$ 与抛物线交于 $N$，正确性质是
    \[
        PN=NM=\frac{PM}{2}.
    \]
    它只说明 $N$ 是\textbf{线段 $PM$ 的中点}，
    不能据此断言抛物线从 $A$ 到 $N$ 与从 $N$ 到 $B$ 的弧长相等。
    \textbf{提醒：}坐标平均与曲线弧长相等是两回事。

    \medskip
    \textbf{\textcolor{CTrap}{易错七：把圆的切线三角形性质直接套用到抛物线}}

    \textbf{圆的情况：}同一圆外一点引出的两条切线段 $PA=PB$，
    故 $\triangle PAB$ 为等腰三角形，底边中线 $PM$ 也是高。

    \textbf{抛物线的情况：}两切线虽然同样交于 $P$，
    但 $PA$、$PB$ 一般\textbf{不相等}。以
    $x^2=4y$ 上 $A(2,1),B(4,4)$ 为例，切线交点为 $P(3,2)$，
    用两点间距离公式：
    \[
        PA=\sqrt{(3-2)^2+(2-1)^2}=\sqrt2,
        \quad PB=\sqrt{(3-4)^2+(2-4)^2}=\sqrt5.
    \]
    因而 $PA\ne PB$。实际可证明的是 $PM$ 竖直，
    \textbf{不是} $PM\perp AB$。
    \textbf{正确面积：}
    \[
        \boxed{S=\frac12 PM\,|x_1-x_2|},
    \]
    这里 $|x_1-x_2|$ 是\textbf{两切点的水平距离}，并不是斜边 $AB$ 的长度。

\end{trapbox}
